% =============================================================================
\section{How the results were reached}
\label{app:provenance}

The work was kept in a running findings ledger and a task log, with corrections recorded as
they happened. The question about tetrahedra was put to a Claude session on 12~October~2026:
does the tetrahedron connect to Jordan algebras? The answer was that it marks where the
exceptional algebra stops. Four orthogonal idempotents force associative coordinates, so
there is no octonionic algebra of rank four, and the session knew of no natural tetrahedron
inside the Albert algebra. Two days later a ChatGPT session exhibited one. Four primitive
idempotents with pairwise overlaps $\frac13$, the overlaps of the qubit tetrahedron of
quantum measurement theory, generate all $27$ dimensions, as an integer determinant of
$-256$ certifies.

Between 17 and 19~October the quantum tetrahedron was worked out over all four rank-three
Jordan algebras, with $\mathrm{SO}(3)$, $\mathrm{SU}(3)$, $\mathrm{Sp}(6)$ and $F_4$ on the
faces, and with the $\mathrm{SU}(2)$ case of Barbieri and of Baez and Barrett reproduced as a
calibration~\cite{Barbieri1998,BaezBarrett1999}. Several tools shared the work. In claude.ai
sessions Claude computed the four rows, their channel decompositions, recoupling matrices and
volume-type spectra, the sixteen-dimensional spin-closure quotient of the exceptional case,
and the bridge between the tetrahedron and the attention layer, in exact rational Python,
python-flint, PARI/GP and LiE built from source. A parallel session checked the exceptional
case again in GAP with a different octonion convention and wrote a thirteen-page draft. GPT
ran long exact-arithmetic computations for two reports, on measurement and spin closure and
on quantum-shape observables in the rank-three Jordan family, with an archive of
certificates. An agent on my own machine extended the $F_4$ sweep to level $140$. The
four-face invariant spaces have dimensions $5$, $8$, $5$ and $5$. For $F_4$, opening one leg
of a four-face invariant gives exactly one of the attention layer's equivariant value maps,
so the shapes of the exceptional tetrahedron are the five channels the layer already used.
One correction from that stage set the habit of keeping a cumulative record. An assembled
``complete'' package turned out to have dropped two of the four earlier write-ups. They had
to be restored from the original archives, and the build scripts showed how the omission had
happened.

The turn to dynamics, described in \cref{sec:history}, led to a second Claude chat, which
opened with a request to prove loop quantum gravity. Claude's first answer was that a
physical theory cannot be proved, and it listed what is open. My next message attached a
statement of a closure target and a criticism, and with reviews I supplied the chat
narrowed to one checkable claim: on Varadarajan's off-shell space~\cite{Varadarajan2022}, can
a correction $T$ make $H-2T$ close with the Lorentzian sign? The chat ended with a no-go
theorem for Lorentzian corrections that preserve the vertices, which an AI referee checked
and which went into my research ledger on 20~October. A Codex session continued the
calculation on the evening of 19~October, and it then moved to Claude Code, with a second Claude
Code session working alongside. The tetrahedron came back twice. On the intertwiner space of
the parent vertex, one candidate correction acted as a real multiple of
$i\,J_1\cdot(J_2\times J_3)$. And a tangent-sensitive four-valent volume carries the sign
factor $\epsilon_{123}-\epsilon_{124}+\epsilon_{134}-\epsilon_{234}$, which equals $\pm4$
exactly at positively spanning vertices. Two days of campaigns found no closure. Several
apparent successes took their sign from conventions the test states cannot see, or diverged
in the ordered limits. An internal referee overturned a headline result, ``memoryless
corrections close with the Euclidean sign'', with a counterexample, and it was retracted.
Early in this work I had set myself the goal of a large new result. A general law whose
constant could be predicted in advance did not count.

The work turned to quantum geometry through the volume. In self-dual variables, with
$\mathrm{SU}(1,1)$ labels, the four-valent volume becomes an infinite zero-diagonal Jacobi
matrix that is not essentially self-adjoint. That observation, the Bohr--Sommerfeld rule it
led to, and their generalisation are the subject of~\cite{FitchLimitCircle}. A search of
the literature for named open problems on which these tools could act then found Bianchi
and Haggard's question~\cite{BH2011}. On 21~October, within a day, came a uniform
four-valent gap; the sharp four-valent law, through the observation that the four-valent
matrix is a hydrogen Stark matrix times a smooth multiplier (\cref{thm:sharp4}); an
explicit five-valent family on which the gap closes, through a dimerised chain; the
six-valent ladder; the Rovelli--Smolin contrast; and the normal form of
\cref{thm:nf-intro}. A referee of the six-valent result found its rate wrong: the exponent
is $n\ln n$, not $\frac32\,n\ln n$. Another check found that a briefing had confused the
normalisations $4T$ and $8T$. When I asked what structural property separates the families
that keep a gap from those that lose it, the answer was already in the normal form: the
ratio $t$ of the two weights.

The equal-weight case took most of the remaining time. The slice $j_1=\frac32$ was first
closed with a floating-point sweep that allowed for rounding. An independent referee found
that this was not rigorous as written and that a tail estimate was missing. The tail lemmas
were then written out, and the finite box was re-run on a GPU in outward-rounded interval
arithmetic over $64\,352\,004$ parameter sets, with smallest certified bound $3.777422$
against the constant $3.7$ of \cref{thm:z32} (\cref{app:z32}). The slices $j_1=1$ and $2$
were proved, and every tuple with spins at most $19$ was certified. One strategy for large
spins was refuted by its own test vector, which became \cref{thm:N1}. Three further rounds
aimed at $j_1\ge\frac52$ (a layer decomposition, wall modes and an $\mathfrak{su}(2)$
ladder) did not succeed, and a fourth was stopped at its start. I ended the campaign on
22~October. That case is \cref{conj:B}, and the ladder picture is in \cref{subsec:ladder}.

The paper was written from the primary sources from 21~October. It was then read by an
independent Claude referee, by Gemini, by a Claude session on claude.ai and by GPT, and each
report was triaged against the source and checked with our own code before anything was
adopted. The Claude referee asked for written proofs of two lemmas of \cref{app:valence6}
that had only been stated, for a full proof of the three-spin pencil
(\cref{thm:pencil}), which had only been sketched, and for a fairer attribution of the
four-valent law to the numerics of Brunnemann and Thiemann~\cite{BT2006}. It also asked us to
remove an unverified assumption about how the runtime library rounds square roots, and the
rigorous-float certificates were re-run with every square root checked in exact arithmetic,
with identical minima on all $2\,274\,052$ tuples. Writing out one of the missing proofs
exposed an inequality slip in the argument it came from, which was repaired with no change
to the theorem. Gemini had read the text layer of the PDF, and nine of its typographical
findings were artefacts of the extraction; its points of wording were adopted. The
claude.ai referee showed that \cref{thm:nf-intro}, as then worded for non-coplanar tangents,
was false, because $1\,680$ sign vectors have a coplanar triple. The theorem is now stated
for tangents in general position, with the remaining case described in
\cref{rem:nongeneric}. The same referee showed that a heuristic limit stated in
\cref{sec:large} was contradicted by larger computations, and it asked for a statement of how
AI was used, which is \cref{sec:history}. GPT found that the reciprocal clause of
\cref{thm:closing} holds only after the two pairs of spins are exchanged. It found that the
tail certificate for $j_1=\frac32$ had entered its decimal constants as binary64 numbers
before enclosing them; re-run with the constants as exact rationals, it passes in all eight
regimes. And it proposed the exact slope $8\nu_0=4\sqrt{10-6\sqrt2}$ for the linear law on
three bounded legs, which \cref{thm:C6} now states and \cref{thm:slope} shows cannot be
improved.
